\chapter{Objective}\label{ch:objective}

\section{The question}
The theory measures the security of a learning procedure in the units of
learning itself. A learning procedure trains models on data held by one or more
parties, and an adversary wants to distinguish, invert or steer. The theory asks
how much an adversary must spend to succeed, and in which currency: samples,
training runs, queries, or the compute of a stated class of learners.

\section{The objective, adopted on 2026-10-25}
Price an adversary against a learning procedure in the same units as the
procedure's own training. Prove the prices by reduction to learning problems
whose thresholds are known, with or without the learning oracle. Measure the
same prices by experiment, the oracle supplying the training runs. The headline
result would be a natural procedure where the attacker's price exceeds the
trainer's by a margin that is named and measured. Halil adopted this objective
on 2026-10-25. A mapping of a security problem onto a hard problem from the
literature is a result in its own right, whatever law or procedure carries
the hardness, and an objection to an instance does not replace a mapping.

\section{Two sides of every question}
The statistical side asks what an adversary can do with unlimited computation
from the information it sees. Its tools are total variation, hockey-stick
divergences, sample complexity of tests and rate distortion. The computational
side asks what an adversary that must learn efficiently can do, for instance a
learner trained by gradient methods. Its tools are statistical-query and
low-degree lower bounds, the exponents that set sharp sample thresholds, and
reductions between learning problems. A result of this theory states which side
it is on, and the most valuable results relate the two.

\section{Constraints set by Halil}
\begin{itemize}[nosep]
\item A cryptographic plug-in placed into a procedure does not count. Finding
  that some aspect of a model reduces to or yields a pseudorandom generator or a
  one-way function would count.
\item Every quantity must be realizable, or approximately realizable, by
  experiment.
\item Polynomial-time adversaries are admitted. Definitions alone are not results.
\item Beautiful mathematics comes first. Any of the three rules above may be broken
  for a result whose mathematics is beautiful, and ugly mathematics is not wanted
  even when it obeys them.
\end{itemize}
